\section{Dashboard Technical Architecture}

\begin{frame}{Technical Architecture - System Design}
  \begin{block}{Component Overview}
    \scriptsize
    The dashboard follows an API-first, non-invasive architecture that
    consumes Keylime's existing REST APIs without requiring any
    modifications to Keylime components.
  \end{block}

  \vspace{0.1cm}

  \begin{center}
    \scriptsize
    \begin{tabular}{|l|l|p{5cm}|}
      \hline
      \textbf{Layer} & \textbf{Component} & \textbf{Technology} \\
      \hline
      \hline
      Frontend & Web UI (SPA) & React.js + TypeScript + D3.js \\
      \hline
      Frontend & WebSocket Client & Real-time updates from backend \\
      \hline
      Backend & API Server & Rust (Axum) + apalis workers \\
      \hline
      Backend & Data Collector & Periodic polling of Keylime APIs \\
      \hline
      Storage & Time-Series DB & TimescaleDB or InfluxDB \\
      \hline
      Storage & Cache & Redis for real-time data \\
      \hline
      Auth & mTLS Client & Admin certificate for Keylime access \\
      \hline
    \end{tabular}
  \end{center}

  \vspace{0.2cm}

  \begin{alertblock}{Key Design Decisions}
    \scriptsize
    \begin{itemize}
      \item \textbf{API-First}: No direct database access -- only public REST APIs
      \item \textbf{Non-Invasive}: Zero modification to Keylime components
      \item \textbf{Secure}: Rust backend + mTLS + rustls (pure-Rust TLS)
      \item \textbf{Scalable}: Stateless frontend, async Rust backend
    \end{itemize}
  \end{alertblock}
\end{frame}

\begin{frame}{Technical Architecture - Data Flow}
  \begin{center}
    \begin{tikzpicture}[scale=0.72, transform shape,
        comp/.style={rectangle, draw, rounded corners, fill=keylimeblue!15, text width=2cm, align=center, minimum height=1cm, font=\tiny\bfseries},
        arrow/.style={->, >=stealth, thick}]

      \node[comp, fill=infoblue!15] (browser) at (0,0) {Browser\\(Web UI)};
      \node[comp, fill=warningyellow!15] (backend) at (4.5,0) {Dashboard\\Backend};
      \node[comp] (verifier) at (9,1) {Keylime\\Verifier};
      \node[comp] (registrar) at (9,-1) {Keylime\\Registrar};

      \draw[arrow] (browser) -- node[above, font=\tiny] {REST/WS} (backend);
      \draw[arrow] (backend) -- node[above, font=\tiny] {mTLS} (verifier);
      \draw[arrow] (backend) -- node[below, font=\tiny] {mTLS} (registrar);
    \end{tikzpicture}
  \end{center}

  \vspace{0.2cm}

  \begin{columns}[T]
    \begin{column}{0.48\textwidth}
      \begin{block}{Frontend Stack (React.js)}
        \scriptsize
        \begin{itemize}
          \setlength{\itemsep}{0pt}
          \item React.js 18+ (SPA)
          \item TypeScript for type safety
          \item D3.js / Recharts for charts
          \item WebSocket real-time updates
          \item Material UI components
          \item React Native path to mobile apps
        \end{itemize}
      \end{block}
    \end{column}
    \begin{column}{0.48\textwidth}
      \begin{block}{Backend Stack (Rust/Axum)}
        \scriptsize
        \begin{itemize}
          \setlength{\itemsep}{0pt}
          \item Axum 0.8 (Tokio-native) API server
          \item apalis for background polling jobs
          \item fred (Redis client, async-first)
          \item SQLx + TimescaleDB for metrics
          \item reqwest + rustls for mTLS client
          \item tracing for structured observability
        \end{itemize}
      \end{block}
    \end{column}
  \end{columns}
\end{frame}

\begin{frame}{Technical Architecture - Why React.js}
  \begin{block}{Frontend Framework Choice: React.js}
    \scriptsize
    React.js is chosen over alternatives (Vue.js, Angular) for its ecosystem maturity,
    component reuse potential, and clear migration path to native mobile applications.
  \end{block}

  \vspace{0.05cm}

  \tiny
  \begin{columns}[T]
    \begin{column}{0.48\textwidth}
      \textbf{React.js Advantages:}
      \begin{itemize}
        \setlength{\itemsep}{0pt}
        \item Largest ecosystem and community
        \item Mature component libraries (MUI, Ant Design)
        \item Strong TypeScript integration
        \item Virtual DOM for efficient re-renders (real-time dashboards)
        \item React Hooks for state management
        \item Wide availability of developers
      \end{itemize}

      \vspace{0.05cm}

      \textbf{Key Libraries:}
      \begin{itemize}
        \setlength{\itemsep}{0pt}
        \item \textbf{Recharts / D3.js}: attestation charts and analytics
        \item \textbf{TanStack Query}: server state caching, polling, refetch
        \item \textbf{React Router}: SPA navigation
        \item \textbf{Zustand / Redux}: client state management
        \item \textbf{Socket.IO client}: WebSocket real-time updates
      \end{itemize}
    \end{column}
    \begin{column}{0.48\textwidth}
      \textbf{Mobile Migration Path (React Native):}
      \begin{itemize}
        \setlength{\itemsep}{0pt}
        \item Shared business logic and TypeScript types between web and mobile
        \item React Native for Android \& iOS from the same codebase
        \item Mobile app use cases: on-call alert notifications (push), fleet health overview, acknowledge/escalate alerts remotely, agent status lookup
        \item Shared API client layer (axios/fetch)
        \item Gradual migration: web-first, mobile later
      \end{itemize}

      \vspace{0.05cm}

      \textbf{Code Sharing Strategy:}
      \begin{itemize}
        \setlength{\itemsep}{0pt}
        \item Monorepo: \texttt{packages/shared}, \texttt{packages/web}, \texttt{packages/mobile}
        \item Shared: API client, types, validation, state logic
        \item Platform-specific: UI components, navigation
      \end{itemize}
    \end{column}
  \end{columns}
\end{frame}

\begin{frame}{Technical Architecture - Why Rust (Backend)}
  \begin{block}{Backend Language Choice: Rust}
    \scriptsize
    Rust is chosen over Python/Go for its memory safety guarantees, ownership model,
    and zero-cost abstractions --- critical for a security-sensitive component that holds
    admin mTLS credentials.
  \end{block}

  \vspace{0.05cm}

  \tiny
  \begin{columns}[T]
    \begin{column}{0.48\textwidth}
      \textbf{Security Advantages:}
      \begin{itemize}
        \setlength{\itemsep}{0pt}
        \item Memory safety without garbage collector
        \item Ownership model prevents data races at compile time
        \item No null pointer dereferences, no buffer overflows
        \item \texttt{rustls}: pure-Rust TLS (no OpenSSL dependency)
        \item \texttt{\#![forbid(unsafe\_code)]} on dashboard crate
        \item Aligns with Keylime's Rust agent direction
      \end{itemize}

      \vspace{0.05cm}

      \textbf{Performance:}
      \begin{itemize}
        \setlength{\itemsep}{0pt}
        \item 400K--585K req/s (TechEmpower benchmarks)
        \item Lowest memory footprint per connection
        \item Ideal for containerized/air-gapped deployments
        \item Native async (Tokio) --- no GIL bottleneck
      \end{itemize}
    \end{column}
    \begin{column}{0.48\textwidth}
      \textbf{Rust Backend Stack:}
      \begin{center}
        \begin{tabular}{|l|l|}
          \hline
          \textbf{Layer} & \textbf{Crate} \\
          \hline
          Web framework & Axum 0.8 \\
          \hline
          Async runtime & Tokio \\
          \hline
          HTTP/mTLS client & reqwest + rustls \\
          \hline
          WebSockets & Axum built-in \\
          \hline
          Background jobs & apalis \\
          \hline
          Database & SQLx (TimescaleDB) \\
          \hline
          Redis & fred \\
          \hline
          Observability & tracing \\
          \hline
        \end{tabular}
      \end{center}

      \vspace{0.05cm}

      \textbf{Ecosystem Coherence:}
      \begin{itemize}
        \setlength{\itemsep}{0pt}
        \item All crates are Tokio-native
        \item Tower middleware shared across web server, background jobs, and HTTP client
        \item Compile-time SQL validation (SQLx)
        \item Single binary deployment (no runtime deps)
      \end{itemize}
    \end{column}
  \end{columns}
\end{frame}

\begin{frame}{Technical Architecture - Event-Driven Ingestion}
  \begin{block}{Pivot from Polling to Events}
    \scriptsize
    Pure HTTP polling cannot scale beyond $\sim$1,000 agents without competing with
    attestation traffic on the Verifier. The dashboard uses an event-driven
    architecture with polling as fallback only.
  \end{block}

  \vspace{0.05cm}

  \tiny
  \begin{columns}[T]
    \begin{column}{0.48\textwidth}
      \textbf{Primary: Event-Driven (Recommended):}
      \begin{itemize}
        \setlength{\itemsep}{0pt}
        \item Keylime ZeroMQ revocation events consumed directly
        \item Propose upstream: state-change webhook/AMQP emit
        \item Message broker (RabbitMQ/Kafka) decouples load
        \item Dashboard subscribes to state-change stream
        \item Near real-time ($<$1s) without polling
        \item Zero additional load on Verifier
      \end{itemize}

      \vspace{0.05cm}

      \textbf{Fallback: Adaptive Polling:}
      \begin{itemize}
        \setlength{\itemsep}{0pt}
        \item Only for environments without event support
        \item List poll: 30s (not 10s) with ETag/If-Modified
        \item Detail poll: on-demand only (user clicks)
        \item Backpressure: double interval if p95 $>$ 500ms
        \item Pause detail polling if p95 $>$ 2s
      \end{itemize}
    \end{column}
    \begin{column}{0.48\textwidth}
      \textbf{Data Sources \& Strategy:}
      \begin{center}
        \begin{tabular}{|l|l|}
          \hline
          \textbf{Data} & \textbf{Ingestion} \\
          \hline
          Agent state & Event stream \\
          \hline
          Attestation result & Event stream \\
          \hline
          Revocations & ZeroMQ events \\
          \hline
          Agent detail & On-demand API \\
          \hline
          IMA policies & Cache (60s TTL) \\
          \hline
          MB policies & Cache (60s TTL) \\
          \hline
          Registrar & Cache (30s TTL) \\
          \hline
        \end{tabular}
      \end{center}

      \vspace{0.05cm}

      \textbf{Split-Brain Mitigation:}
      \begin{itemize}
        \setlength{\itemsep}{0pt}
        \item Event sequence numbers detect gaps
        \item Periodic reconciliation sweep (every 5 min)
        \item Dashboard marks data as ``stale'' if gap detected
        \item UI shows last-verified timestamp per agent
      \end{itemize}
    \end{column}
  \end{columns}
\end{frame}

\begin{frame}{Technical Architecture - IMA Log \& Data Decoupling}
  \begin{block}{Avoiding IMA Log DDoS on the Verifier}
    \scriptsize
    IMA measurement logs can be megabytes per agent. Proxying them through the
    dashboard backend and Verifier under concurrent SOC analyst load will crash the Verifier.
    Logs must be decoupled from the real-time monitoring path.
  \end{block}

  \vspace{0.1cm}

  \scriptsize
  \begin{columns}[T]
    \begin{column}{0.48\textwidth}
      \textbf{Log Decoupling Architecture:}
      \begin{itemize}
        \setlength{\itemsep}{0pt}
        \item Keylime pushes IMA/boot logs to centralized log store (Elasticsearch, Splunk, Loki)
        \item Dashboard queries SIEM directly for historical logs
        \item Never proxy bulk logs through Verifier API
        \item Rate-limit concurrent log fetches (max 5 parallel)
        \item Streaming pagination for large logs
      \end{itemize}

      \vspace{0.1cm}

      \textbf{Concurrency Protection:}
      \begin{itemize}
        \setlength{\itemsep}{0pt}
        \item Per-user request rate limiting
        \item Global concurrent API call cap
        \item Circuit breaker on Verifier latency
        \item Queue overflow: reject with 429
      \end{itemize}
    \end{column}
    \begin{column}{0.48\textwidth}
      \textbf{What the Dashboard Caches:}
      \begin{itemize}
        \setlength{\itemsep}{0pt}
        \item Attestation \textit{results} (pass/fail + timestamp)
        \item Agent metadata (UUID, IP, state)
        \item Policy checksums and assignment
        \item Certificate expiry dates
      \end{itemize}

      \vspace{0.1cm}

      \textbf{What is NEVER Cached:}
      \begin{itemize}
        \setlength{\itemsep}{0pt}
        \item Raw TPM quotes
        \item IMA measurement logs
        \item Boot event logs
        \item PoP tokens or session secrets
      \end{itemize}

      \vspace{0.1cm}

      \textbf{API Version Compatibility:}
      \begin{itemize}
        \setlength{\itemsep}{0pt}
        \item Negotiate API version on startup (\texttt{GET /versions})
        \item Support v2 and v3 simultaneously
        \item Graceful degradation on unknown fields
        \item Schema migration tool for dashboard upgrades
      \end{itemize}
    \end{column}
  \end{columns}
\end{frame}

\begin{frame}[fragile]{Technical Architecture - Deployment \& API Access}
  \begin{block}{Deployment Options}
    \scriptsize
    The dashboard can be deployed through multiple mechanisms to fit different
    infrastructure environments.
  \end{block}

  \vspace{0.1cm}

  \scriptsize
  \begin{center}
    \begin{tabular}{|l|p{4cm}|p{3.5cm}|}
      \hline
      \textbf{Method} & \textbf{Description} & \textbf{Best For} \\
      \hline
      \hline
      Container (Podman/Docker) & Pre-built OCI image & Quick deployments, CI/CD \\
      \hline
      Kubernetes (Helm chart) & Helm-managed deployment & Production clusters \\
      \hline
      RPM package & Native RHEL/Fedora package & Bare-metal, VM installs \\
      \hline
      Systemd service & Direct service management & Single-node setups \\
      \hline
    \end{tabular}
  \end{center}

  \vspace{0.2cm}

  \begin{exampleblock}{Keylime API Access Example}
    \begin{lstlisting}[language=bash, numbers=none, basicstyle=\ttfamily\tiny]
# List all agents from verifier (mTLS required)
curl --cert admin.pem --key admin-key.pem \
  https://verifier:8881/v2/agents/

# Get specific agent status
curl --cert admin.pem --key admin-key.pem \
  https://verifier:8881/v2/agents/{agent_id}
    \end{lstlisting}
  \end{exampleblock}
\end{frame}
